\begin{exercise}{93}
Prove or disprove the following statements:
\begin{enumerate}[label=(\arabic*)]
\item For every $\infty$-category $\mathcal C$, a set $S$ of morphisms of $\mathcal C$
and a set $T$ of morphisms of $\mathcal C[S^{-1}]$, the functor
$\mathcal C\to\mathcal C[S^{-1}][T^{-1}]$ is a localization of $\mathcal C$.
\item For every $\infty$-category $\mathcal C$, a set $S$ of morphisms of $\mathcal C$
and a set $T$ of morphisms of $\mathcal C$, the functor
$\mathcal C\to\mathcal C[S^{-1}][T^{-1}]$ is a localization of $\mathcal C$.
Here, we view morphisms of $\mathcal C$ as morphisms of $\mathcal C[S^{-1}]$
via the canonical functor $\mathcal C\to\mathcal C[S^{-1}]$.
\end{enumerate}

\begin{proof}
(2) This is true: the composite localizes along $S\cup T$.
Write the two localization functors as $p\colon\mathcal C\to\mathcal M$
and $r\colon\mathcal M\to\mathcal Q$.
For every $\infty$-category $\mathcal E$, the restriction
$(rp)^*=p^*r^*$ is fully faithful, since both factors are
(Definition~2.4.2).
Its essential image consists exactly of the functors inverting $S\cup T$:
such a functor factors through $p$ up to natural equivalence,
and the factor inverts $p(T)$, hence factors through $r$.
Conversely, every functor factoring through $rp$ inverts both sets.
Here natural equivalences preserve the set of inverted morphisms,
as proved in Exercise~\ref{ex:92}.
This verifies Definition~2.4.2 for the composite.

(1) This is false.
Let $\mathcal C$ be the nerve of the path category of
\[
 A\xrightarrow{k}V\xleftarrow{s}U\xrightarrow{h}B.
\]
Its only nonidentity morphisms are $k,s,h$.
Choose a localization $p\colon\mathcal C\to\mathcal M$ along $\{s\}$.
By Corollary~1.2.13, choose an edge $t$ of $\mathcal M$ representing
\[
 [t]=[p(h)][p(s)]^{-1}[p(k)]\quad\text{in }h\mathcal M.
\]
Let $r\colon\mathcal M\to\mathcal Q$ localize along $\{t\}$, and put $q=rp$.
These localizations exist by Proposition~2.4.8.

Let $\mathcal E$ be the full subcategory of sets on
$1=\{0\}$ and $2=\{0,1\}$; write $i\colon1\to2$ for the inclusion
and $!\colon2\to1$ for the unique map.
Define two functors from the path category to $\mathcal E$ by
\[
 \begin{array}{c|cccc|ccc}
       & A & U & V & B & k & s & h \\
 F     & 1 & 2 & 2 & 1 & i & \mathrm{id}_2 & ! \\
 G     & 2 & 1 & 1 & 2 & ! & \mathrm{id}_1 & i
 \end{array}
\]
and also denote their nerve maps by $F,G$.
For any $H\colon\mathcal M\to\mathrm N(\mathcal E)$ with
$Hp\simeq F$, the image of $[t]$ under $hH$ is, up to isomorphisms
at its endpoints,
$F(h)F(s)^{-1}F(k)=!i=\mathrm{id}_1$.
Indeed, the induced natural isomorphism on homotopy categories from
Exercise~\ref{ex:89} conjugates each of the three factors;
the intermediate isomorphisms cancel.
Thus the factor of $F$ through $p$ inverts $t$ and extends through $r$.
Since $F(k)=i$ and $F(h)=!$ are not isomorphisms, neither $q(k)$ nor
$q(h)$ is an equivalence.
Consequently the morphisms of $\mathcal C$ inverted by $q$ are exactly
the identities and $s$.

If $q$ were a localization along some set $R$ of morphisms of $\mathcal C$,
then $R$ would be contained in this set.
The functor $G$ inverts all of it, so it would factor as $Zq\simeq G$.
Apply the same cancellation calculation to $Zr$:
the image of $[t]$ would be, up to endpoint isomorphisms,
$G(h)G(s)^{-1}G(k)=i!$, a nonbijective constant map $2\to2$.
But $r(t)$ is an equivalence, so its image under $Z$ must be an
equivalence. This contradiction proves that $q$ is not a localization
along any set of morphisms of $\mathcal C$.
\end{proof}
\end{exercise}
